% Build:
% cd /home/tian/workspace/dpol/smsimulator5.5/docs/reports/reconstruction
% latexmk -pdf fdc_reconstruction_detailed.tex

\documentclass[aspectratio=169]{beamer}
\usetheme{Madrid}
\usecolortheme{default}

\usepackage{amsmath}
\usepackage{booktabs}
\usepackage{siunitx}

\title[FDC Reconstruction in ANAROOT]{Detailed FDC Charged-Particle Reconstruction in ANAROOT/SAMURAI}
\subtitle{From Drift Hits to Target Momentum}
\author{SMSimulator5.5 / ANAROOT Notes}
\date{\today}

\begin{document}

\begin{frame}
\titlepage
\end{frame}

\begin{frame}{Scope and Goal}
\textbf{Goal}
\begin{itemize}
\item Explain the FDC-based charged-track reconstruction chain implemented in ANAROOT.
\item Clarify how target momentum is reconstructed and where fallback methods are used.
\end{itemize}

\vspace{0.2cm}
\textbf{Covered modules}
\begin{itemize}
\item Hit calibration and drift-length conversion
\item Local FDC track fit (1D/3D)
\item Global Runge-Kutta fit in field map
\item MultiDimFit and matrix fallback branches
\end{itemize}
\end{frame}

\begin{frame}{Top-Level Reconstruction Branches}
Inside \texttt{TArtRecoFragment::ReconstructData}, reconstruction uses:
\begin{enumerate}
\item \textbf{Runge-Kutta global fit} (preferred when field map and RK config are ready)
\item \textbf{MultiDimFit polynomial regression} (if pre-generated fit files are ready)
\item \textbf{0th/1st-order matrix fallback} (if above are not available)
\end{enumerate}

\vspace{0.2cm}
The charged fragment object \texttt{TArtFragment} is filled by one of these branches.
\end{frame}

\begin{frame}{FDC Hit Input and Drift Length Conversion}
\textbf{Input containers}
\begin{itemize}
\item FDC raw hits are loaded as \texttt{TArtDCHit} objects.
\item Hits are grouped by layer and wire direction (X/U/V/Y) in buffers.
\end{itemize}

\textbf{Drift length calculation logic}
\begin{itemize}
\item If STC function exists: \(dl = f_{\mathrm{stc}}(\mathrm{TDC})\)
\item Else if TDC distribution exists: lookup \(\mathrm{TDC} \to dl\) table
\item Else fallback with constant drift velocity:
\[
  dl = v_{\mathrm{drift}}\,(\mathrm{TDC}_0 - \mathrm{TDC})
\]
\item Unphysical values are rejected (negative / too large drift length).
\end{itemize}
\end{frame}

\begin{frame}{Pattern Building and Left/Right Ambiguity}
Each wire hit has left/right ambiguity around the sense wire.

\vspace{0.2cm}
For a pattern with \(n\) hits in one projection, ANAROOT scans:
\[
2^n \text{ sign combinations}
\]
where each hit position candidate is
\[
x_i^{\pm} = x_{\mathrm{wire},i} \pm dl_i.
\]

\vspace{0.2cm}
A straight-line hypothesis is fitted for each pattern, and the minimum-\(\chi^2\) candidate is selected.
\end{frame}

\begin{frame}{1D Track Fit in One Projection}
Model in one view:
\[
x(z) = x_0 + a z
\]
with least-squares solve from normal equations.

\vspace{0.2cm}
Objective for each sign pattern:
\[
\chi^2 = \sum_i \left(x_0 + a z_i - x_i\right)^2.
\]

\vspace{0.2cm}
The best candidate stores:
\begin{itemize}
\item fitted intercept \(x_0\), slope \(a\)
\item selected sign-resolved hit coordinates
\item drift lengths, layer IDs, plane IDs
\end{itemize}
\end{frame}

\begin{frame}{3D Combination of X/U/V Tracks}
After best 1D tracks are found for X/U/V projections:
\begin{itemize}
\item ANAROOT solves a 4-parameter linear system for 3D track state.
\item U/V layers are projected with detector stereo angles (typically \(\pm30^\circ\)).
\item Combined \(\chi^2\) is accumulated from residuals in all used views.
\end{itemize}

\vspace{0.2cm}
Output is a \texttt{TArtDCTrack} with position/angle and associated hit-level diagnostics.
\end{frame}

\begin{frame}{Track Candidate Selection in FDC1/FDC2}
\textbf{Within each chamber}
\begin{itemize}
\item Candidate patterns are generated from hit combinations.
\item Candidates with valid \(\chi^2\) and required hit-layer conditions are kept.
\item Tracks are sorted by quality; reconstruction uses the lowest-\(\chi^2\) candidate.
\end{itemize}

\vspace{0.2cm}
\textbf{In fragment reconstruction}
\begin{itemize}
\item One charged particle is assumed.
\item The first entry in FDC1/FDC2 track arrays (best \(\chi^2\)) is taken.
\end{itemize}
\end{frame}

\begin{frame}{Global RK Fit: Initialization}
From FDC2 local track, ANAROOT builds an initial global state near HOD:
\begin{itemize}
\item Position extrapolation: \(x_{\mathrm{HOD}} = x_{\mathrm{FDC2}} + L\,a_x\), similarly for \(y\)
\item Direction from local angles \((u,v)\)
\item Initial momentum magnitude from configured \(B\rho\) center (scaled in code)
\end{itemize}

\vspace{0.2cm}
Runge-Kutta propagation is then performed through the SAMURAI field map across detector planes.
\end{frame}

\begin{frame}{Global RK Fit: Objective Function}
For each hit layer, predicted local hit position is compared with measured one.

\vspace{0.2cm}
Weighted objective:
\[
\chi^2 = \frac{1}{N_{\mathrm{hit}}-5}\sum_i
\frac{\left(x_i^{\mathrm{hit}}-x_i^{\mathrm{calc}}\right)^2}{\sigma_i^2}
\]
where \(\sigma_i\) comes from detector resolution tables.

\vspace{0.2cm}
Five fitted parameters correspond to trajectory state components including \(q\propto 1/p\).
\end{frame}

\begin{frame}{Global RK Fit: Parameter Update Strategy}
ANAROOT uses a Levenberg-Marquardt style step with SVD-based linear solve.

\vspace{0.2cm}
Per iteration:
\begin{enumerate}
\item Build derivatives/Hessian-like matrix from RK hit-point coefficients.
\item Add damping term (LM-like regularization).
\item Solve update vector with SVD (stable under near-singular conditions).
\item Update state and repeat until convergence/stop conditions.
\end{enumerate}

\vspace{0.2cm}
This is a deterministic nonlinear least-squares RK fit, not a Kalman filter.
\end{frame}

\begin{frame}{Outputs at Target and TOF Path Length}
After convergence, \texttt{saveTrackParameters()} stores:
\begin{itemize}
\item Primary position at target plane
\item Primary momentum at target plane
\item Path length target-to-HOD and total traced length
\end{itemize}

These are copied into \texttt{TArtFragment} and later used with HOD/TZero timing to compute:
\[
\beta,\quad \gamma,\quad A/Q.
\]
\end{frame}

\begin{frame}{MultiDimFit Branch (Fast Surrogate)}
If RK branch is not used (or in configured modes), MultiDimFit can provide:
\begin{itemize}
\item rigidity \(B\rho\)
\item path length
\end{itemize}

\vspace{0.2cm}
Input vector in code is 6D:
\[
x = (\text{FDC1 pos/angle terms with BDC-target interpolation},\ \text{FDC2 terms}).
\]

\vspace{0.2cm}
Prediction is a polynomial expansion generated offline by ROOT \texttt{TMultiDimFit::MakeRealCode}.
\end{frame}

\begin{frame}{0th/1st-Order Matrix Fallback}
When neither RK nor MultiDimFit is ready:
\begin{itemize}
\item Code loads 0th and 1st order optics matrices from text files.
\item FDC observables are inserted into linearized transport relation.
\item Inverted sub-matrix is used to solve \(\delta\)-like quantity.
\end{itemize}

\vspace{0.2cm}
This is fast and robust as a fallback, but less expressive than full RK tracing.
\end{frame}

\begin{frame}{Parameter Files and Provenance}
\textbf{MultiDimFit files}
\begin{itemize}
\item Environment variables:
  \texttt{SAMURAI\_MULTIDIM\_FILE\_RIG}, \texttt{SAMURAI\_MULTIDIM\_FILE\_LEN}
\item Typical files: \texttt{db/samurai\_fun\_rig.C}, \texttt{db/samurai\_fun\_len.C}
\item Generated offline by ROOT MultiDimFit from Geant4 simulation trees.
\end{itemize}

\textbf{Matrix fallback files}
\begin{itemize}
\item Environment variables:
  \texttt{SAMURAI\_MATRIX0TH\_FILE}, \texttt{SAMURAI\_MATRIX1ST\_FILE}
\item Usually provided by an external Geant4 matrix calculation workflow.
\end{itemize}
\end{frame}

\begin{frame}{Failure Modes and Debug Checklist}
\textbf{Typical failure causes}
\begin{itemize}
\item Missing FDC1/FDC2 tracks or insufficient hit layers
\item Drift-length calibration out of range (TDC window / pitch constraints)
\item Missing environment-variable parameter files
\item Field-map or detector-configuration mismatch
\end{itemize}

\vspace{0.2cm}
\textbf{Fast checks}
\begin{itemize}
\item Track counts and \(\chi^2\) of FDC1/FDC2 best tracks
\item RK status flag, iteration count, residual array
\item Loaded file paths for MultiDimFit/matrix env vars
\end{itemize}
\end{frame}

\begin{frame}{Key Source Files (Code Reading Map)}
\small
\begin{itemize}
\item \texttt{src/sources/Reconstruction/SAMURAI/src/TArtCalibFDC1Hit.cc}
\item \texttt{src/sources/Reconstruction/SAMURAI/src/TArtCalibDCTrack.cc}
\item \texttt{src/sources/Reconstruction/SAMURAI/src/TArtRecoFragment.cc}
\item \texttt{src/sources/Reconstruction/SAMURAI/src/TArtCalcGlobalTrack.cc}
\item \texttt{src/sources/Reconstruction/SAMURAI/src/RungeKuttaUtilities.cc}
\item \texttt{src/sources/Reconstruction/SAMURAI/src/FieldMan.cc}
\item \texttt{src/sources/Reconstruction/SAMURAI/src/SAMURAIFieldMap.cc}
\end{itemize}

\vspace{0.2cm}
This map is sufficient to trace FDC reconstruction from raw hit to target momentum.
\end{frame}

\begin{frame}{Takeaway}
\begin{itemize}
\item FDC local reconstruction solves wire ambiguity + straight-line fit in projections.
\item Global charged-particle momentum is obtained by RK field-map fit with weighted residual minimization.
\item MultiDimFit and first-order matrix methods are practical alternatives for speed or robustness.
\end{itemize}
\end{frame}

\end{document}
